\documentclass[aspectratio=169,11pt]{beamer}
\usepackage[T1]{fontenc}
\usepackage{lmodern,amsmath,amssymb}
\definecolor{ink}{HTML}{18364D}
\setbeamercolor{structure}{fg=ink}
\setbeamercolor{frametitle}{fg=ink}
\setbeamertemplate{navigation symbols}{}
\setbeamertemplate{footline}{\hfill\insertframenumber\kern1.5em\vskip8pt}
\setbeamersize{text margin left=10mm,text margin right=10mm}
\newcommand{\cI}{\mathcal I}
\newcommand{\cO}{\mathcal O}
\newcommand{\db}{\bar\partial}
\newcommand{\norm}[1]{\lVert#1\rVert}
\title{Curvature and vanishing theorems}
\subtitle{Bochner--Kodaira, Kodaira, Nadel and Kawamata--Viehweg}
\author{Daniel}
\date{Hodge theory}
\begin{document}
\begin{frame}
\titlepage
\end{frame}

% Verbitsky 2025, lecture 16, slide 5.
\begin{frame}{Bochner--Kodaira identity and proof}
\small
On a Kähler manifold, let \(B\) be a Hermitian holomorphic line bundle,
\(\nabla=P+Q=\partial_B+\db_B\), and \(\Lambda=(\omega\wedge)^*\).
\[
\boxed{\Delta_Q=\Delta_P+[i\Theta_B\wedge,\Lambda].}
\]
\textbf{Proof.} The Kähler identities give
\(Q^*=-i[\Lambda,P]\), \(P^*=i[\Lambda,Q]\). Hence
\[
\begin{aligned}
\Delta_Q-\Delta_P
&=-i\bigl(Q[\Lambda,P]+[\Lambda,P]Q
             +P[\Lambda,Q]+[\Lambda,Q]P\bigr)\\
&=-i\bigl(\Lambda(PQ+QP)-(PQ+QP)\Lambda\bigr)\\
&=[i\Theta_B\wedge,\Lambda].
\end{aligned}
\]
The middle equality expands and cancels the \(Q\Lambda P\) and
\(P\Lambda Q\) terms. The last uses
\(P^2=Q^2=0\) and \(PQ+QP=\nabla^2=\Theta_B\wedge\).
\medskip

\footnotesize
Earlier inputs: Kähler identities (lecture 10), Chern connection
(lecture 14). This identity was proved in lecture 16.
\end{frame}

% Verbitsky 2025, lecture 16; HP complex-geometry.tex.
\begin{frame}{Kodaira vanishing and proof}
\small
\textbf{Theorem.} Let \((X,\omega)\) be compact Kähler, \(\dim_{\mathbb C}X=n\),
and \(B\) have a smooth metric with \(i\Theta_B>0\). Then
\[
\boxed{H^q(X,K_X\otimes B)=0\qquad(q>0).}
\]
\textbf{Proof.} Compactness gives \(i\Theta_B\ge\varepsilon\omega\).
Represent a class by a harmonic \(B\)-valued \((n,q)\)-form \(u\).
In a unitary frame diagonalizing curvature, with eigenvalues
\(\lambda_j\ge\varepsilon\), the curvature operator on the monomial
indexed by \(J\), \(|J|=q\), has eigenvalue
\(\sum_{j\in J}\lambda_j\ge q\varepsilon\).
Thus Bochner--Kodaira gives
\[
0=\langle\Delta_{\db}u,u\rangle
=\norm{\partial_Bu}^2+\norm{\partial_B^*u}^2
+\langle[i\Theta_B\wedge,\Lambda]u,u\rangle
\ge q\varepsilon\norm{u}^2.
\]
For \(q>0\), \(u=0\). Every class therefore vanishes.
\medskip

\footnotesize
Earlier inputs: Dolbeault resolution (lectures 12--13), bundle-valued
Hodge theory (15), curvature-operator calculation (16).
\end{frame}

% Verbitsky 2025, lecture 23, slides 5--12.
\begin{frame}{Nadel vanishing and proof}
\small
\textbf{Theorem.} Let \((X,\omega)\) be compact Kähler and \(B\) have a
singular metric \(h\) with \(i\Theta_h\ge\varepsilon\omega\), \(\varepsilon>0\),
as currents. Then
\[
\boxed{H^q(X,K_X\otimes B\otimes\cI_h)=0\qquad(q>0).}
\]
Here \((\cI_h)_x=\{f\in\cO_{X,x}:|f|^2e^{-2\varphi}\in L^1_{\rm loc}\}\)
when \(|e|_h^2=e^{-2\varphi}\).
\medskip

\textbf{Proof, using weighted \(L^2\) solvability.}
The local weighted Dolbeault lemma gives a fine resolution of
\(K_X\otimes B\otimes\cI_h\) by locally \(L^2_h\), \(B\)-valued
\((n,q)\)-forms with weak \(\db\) in \(L^2_h\).
A class has a global closed representative \(f\).
The singular \(L^2\) existence theorem gives
\[
\db u=f,\qquad \norm{u}_h^2\le(q\varepsilon)^{-1}\norm{f}_h^2.
\]
Thus every class is exact.
\medskip

\footnotesize
Earlier tools: closed operators (17), complete-metric \(L^2\) theory (18),
psh exhaustion (21), singular psh approximation (22).
The weighted lemma and singular solvability are the analytic inputs
to establish in lecture 23.
\end{frame}

% Verbitsky 2025, lecture 24, slides 4, 13--16.
\begin{frame}{Kawamata--Viehweg vanishing and proof}
\small
\textbf{Theorem.} On a smooth projective \(n\)-fold \(X\), if \(L\) is
nef and big, then \(H^q(X,K_X\otimes L)=0\) for \(q>0\).
Nef means \(L\cdot C\ge0\) for every curve; big means
\(\limsup_m h^0(X,L^m)/m^n>0\).
\medskip

\textbf{Kawamata--Viehweg:} obtain vanishing under the weaker
positivity assumptions \textquotedblleft nef and big.\textquotedblright
\medskip

\textbf{Proof.}
Choose a smooth very ample divisor \(A\).
The restriction sequence and
\(h^0(A,L^m|_A)=O(m^{n-1})\) give
\(L^d=\cO_X(A+D)\) with \(D\) effective.
For \(N>d\),
\[
L^N=(L^{N-d}\otimes\cO_X(A))\otimes\cO_X(D).
\]
The first factor is ample by nefness.
Its positive metric times the divisor metric, followed by the
\(N\)-th root, gives strictly positive curvature on \(L\).
The weight is locally a smooth factor times \(|f_D|^{-2/N}\).
For fixed \(D\) and large \(N\), it is integrable, so \(\cI_h=\cO_X\).
Nadel proves the claim.
\medskip

\footnotesize
Inputs: Poincaré--Lelong (22); small-exponent integrability and section
growth (24); nef plus ample is ample; positive metrics on ample bundles.
\end{frame}
\end{document}
